\oldpage[IV-2]{177}
Sean $K$ un cuerpo algebraicamente cerrado de característica 0
y $B$ el dominio entero
$K[U,V,W]/(U^2(U\setmin W)\setmin V^2(U+W))$
($U$, $V$ y $W$ indeterminadas), de modo que
$Y=\operatorname{Spec}(B)$ es un «cono con vértice en el origen y
una generatriz doble». Denotaremos por $u$, $v$ y $w$ las imágenes
de $U$, $V$ y $W$ en $B$. Sea $R$ el cuerpo de fracciones de $B$
y consideremos en $R$ el elemento $t=v(u+w)/u$, que no pertenece
a $B$; demostremos que $C=B[t]$ es la \emph{clausura integral} de $B$.
En efecto, $t^2=u^2\setmin w^2$, de modo que $t$ es entero sobre $B$,
y $v=tu/(u+w)$. El anillo $C_1=K[t,u,w]$ es integralmente cerrado,
pues es isomorfo a
$K[T,U,W]/(T^2\setmin U^2+W^2)$ y es, por tanto, la clausura integral
del anillo integralmente cerrado $K[U,W]$ en la extensión
cuadrática $K(U,W)(\sqrt{U^2-W^2})$ de su cuerpo de fracciones
(Bourbaki, \emph{Alg. comm.}, cap. V, §~1, n\textsuperscript{o}~6, prop. 18).
El anillo de fracciones $K[t,u,w,1/(u+w)]$ de $C_1$ es,
pues, también integralmente cerrado. Pongamos $C_2=K[t,v,w]$ y $a=t-v$.
Las identidades anteriores dan
$t(t\setmin v)^2\setmin w^2(2v\setmin t)=0$, es decir, $t(a^2-w^2)=-2aw^2$.
En $(C_2)_a$, la identidad $(1-t/(2a))(a^2-w^2)=a^2$ muestra que
$a^2-w^2$ es invertible. Así pues,
\[
(C_2)_a=K[a,w,1/a,1/(a^2-w^2)],
\]
pues $t=-2aw^2/(a^2-w^2)$ y $v=t-a$. Como además $u=vw/a$,
se tiene $K(a,w)=R$; por tanto, $a$ y $w$ son algebraicamente independientes
sobre $K$. Este localizado de un anillo de polinomios es integralmente
cerrado. En particular,
$K[t,v,w,1/(t\setmin v)]=K[t,\frac{tu}{u+w},w,\frac{u+w}{tw}]$ es integralmente cerrado.\footnote{Corrección editorial. El texto impreso omite el cuadrado de $t-v$ en la ecuación y afirma que $C_2$ es integralmente cerrado. El argumento corregido establece directamente la propiedad para el localizado que se utiliza aquí.} Por último, teniendo en cuenta que $u$
y $w$ son algebraicamente independientes sobre $K$, se prueba
fácilmente que
$C=K[t,u,v,w]=K[t,u,w,1/(u+w)]\cap K[t,v,w,1/(t\setmin v)]$,
lo que termina de demostrar que $C$ es la clausura integral de
$B$. Es inmediato que, si $\mathfrak m_0$ es el ideal maximal de
$B$ engendrado por $u$, $v$ y $w$ («vértice del cono»), existe
un único ideal maximal $\mathfrak n_0$ de $C$ situado sobre
$\mathfrak m_0$, a saber, el ideal engendrado por $t$, $u$, $v$
y $w$. Si se pone $A=B_{\mathfrak m_0}$, se deduce fácilmente
que $A'=C_{\mathfrak n_0}$ es la clausura integral de $A$,
que es, por tanto, unirramoso y también geométricamente
unirramoso, puesto que su cuerpo residual $K$ es algebraicamente
cerrado. Pero en $A$ el ideal primo $\mathfrak p$ engendrado por
$u$ y $v$ es tal que la clausura integral $A_\mathfrak p[t]$
de $A_\mathfrak p$ no es un anillo local.

Veremos, no obstante, más adelante (9.7.10) que, cuando $X$ es
un preesquema localmente noetheriano tal que, si $X'$ es el
normalizado de $X_{\mathrm{red}}$, el morfismo canónico
$X'\to X$ es \emph{finito} (lo que ocurrirá si los anillos de los
abiertos afines de $X$ son \emph{universalmente japoneses} (\textbf 0, 23.1.1)),
el conjunto de los puntos $x\in X$ en los que $X$ es
geométricamente unirramoso es \emph{localmente construible}.
